\section{Next Paper}
\label{sec:nextpaper}
The Checker Framework manual advises developers to ``rewrite your code to be simpler for the checker to analyze; this is likely to make it easier for people to understand, too''~\cite{checkerframework-manual}. Similarly, the documentation of the OpenJML verification tool notes that ``success in checking the consistency of the specifications and the code will depend on the complexity and style in which the code and specifications are written''~\cite{openjml}. These statements reflect a widely held assumption in the software verification community that code that is easier for verification tools to analyze is also easier for humans to understand.

Although this relationship has been widely assumed, prior empirical evidence remains limited in its ability to characterize it under controlled conditions. Recent work~\cite{FeldmanKC2023} investigated this hypothesis by defining verifiability as the effort required to establish safety properties (e.g., the absence of null pointer dereferences or out-of-bounds accesses) using verification tools, and operationalizing it via the presence or absence of false positive warnings. That work conducted a meta-analysis over code snippets from prior comprehensibility studies, aggregating results across different tools, warning types, and experimental settings, and found associations between verification-related signals and measures of human understanding.

However, because the analysis aggregates results from studies that were not designed to isolate this relationship, it cannot cleanly attribute differences in comprehensibility to differences in verifiability. For example, the snippets are drawn from prior studies, so the authors do not control for other aspects of the code, such as its source, structure, and size, which may independently affect comprehension and verifiability. Moreover, these studies rely on a wide range of proxies to measure comprehensibility, many of which we show in our recent work (Section~\ref{ase}) to be weak or unreliable. As a result, the meta-analysis conflates variation in verifiability with variation in both program characteristics and measurement quality, limiting its ability to provide a clear and controlled understanding of the relationship between verifiability and comprehensibility.

To better understand this relationship, we propose a controlled experimental study that systematically manipulates verifiability while constraining other sources of variation. We fix the verification setting by restricting the study to a small, predefined set of verification tools selected based on expert judgment. We construct a dataset of code snippets drawn from high-quality, real-world repositories, selected to satisfy predefined criteria such as moderate size, use of standard libraries, and absence of annotations, in order to reduce variability unrelated to the treatment. Within this setting, each snippet is represented by two semantically equivalent versions: a high-verifiability version that produces no warnings under the selected tools, and a low-verifiability version that produces false positive warnings under the same tools (an example is shown in Figure~\ref{fig:verifiability-example}). These versions are created using carefully designed, semantics-preserving transformations that induce or eliminate false positive warnings, allowing us to manipulate verifiability while holding program functionality constant.

To further control for confounding factors, we select transformations that minimally affect code length and control flow, and mitigate participant- and task-level variability through randomized assignment of snippet versions and appropriate statistical modeling.

Participants complete all study sessions in a controlled environment. They then perform input–output comprehension tasks on the snippets in a within-subject manner, where each participant sees one version of each snippet. For each snippet, they are given a concrete input and asked to determine the corresponding output of the method. Based on the results of our prior work (Section~\ref{ase}), we use the time required to answer these questions as a validated proxy for code comprehensibility. By collecting response times across multiple snippets and participants and analyzing the data using appropriate statistical models, we examine how differences in verifiability are reflected in measures of human understanding.

Our main research question is:

\begin{researchquestions}
\item Are semantically equivalent code snippets with higher verifiability more comprehensible than those with lower verifiability?
\end{researchquestions}
  
If a substantial relationship between verifiability and code comprehensibility is established, verifiability could serve as a practical automatic proxy, reducing reliance on costly human studies and enabling scalable tools for improving code quality. Even a moderate relationship would support incorporating verifiability as a component in multi-factor proxies. Such a proxy could be used to identify code that is difficult to understand and guide refactoring toward more comprehensible forms. This is particularly important given that prior studies estimate that developers spend between 58\% and 70\% of their time understanding code~\cite{Ko2006,Siegmund2014}, suggesting that improvements in code comprehensibility could have a significant impact on developer productivity. Moreover, with the increasing use of large language models for code generation, such proxies could help ensure that generated code remains understandable to human developers.

  
\subsection{Background}
\subsubsection{Code Verifiability}
\label{sec:next-verifiability}

Similar to prior work~\cite{FeldmanKC2023}, we define code verifiability as the effort required for a developer to use a verification tool to establish the correctness of a code snippet with respect to a set of safety properties. In practice, verifiability is defined relative to a specific verification setting, including the tool, the properties being checked, and the available specifications.

Directly measuring verification effort would require controlled studies with developers actively using verification tools, which is difficult to scale. Therefore, we operationalize verifiability using a proxy based on the behavior of sound static analysis tools. Such tools attempt to construct proofs of correctness and may issue warnings when they cannot establish a proof.

We restrict the verification ecosystem to the Checker Framework and select a small subset of its checkers. The Checker Framework is open-source, actively maintained, and widely used for verification in Java. Within this framework, we select checkers based on expert judgment, prioritizing those that are widely used and that produce meaningful safety warnings. For example, the nullness checker is appropriate for our study because it can produce false warnings at any point where a dereference occurs. We determine the final set of checkers in consultation with experts in the Checker Framework, and we also have direct access to a domain expert within our team, which enables more controlled study design and more reliable identification of false positive warnings.

We operationalize verifiability in terms of the presence or absence of false positive warnings produced by the selected checkers. A false positive warning occurs when the tool reports a potential violation in code that is in fact correct, typically due to limitations in the analysis or implicit assumptions in the code. Such warnings require additional developer effort to resolve.

To enable controlled manipulation, we define verifiability as a binary condition:
\begin{itemize}
    \item high verifiability corresponds to code that produces no warnings,
    \item low verifiability corresponds to semantically equivalent code that produces false positive warnings.
\end{itemize}

These choices define the scope of the study. Accordingly, our results should be interpreted as applying to verifiability within the selected Checker Framework setting, rather than to all possible verification tools or warning types.

This operationalization allows us to construct paired code snippets that differ only in whether the checker can successfully verify them.

\subsubsection{Code Comprehensibility}

In prior work~\ref{ase}, we defined code comprehensibility operationally through expert agreement on the relative difficulty of understanding code snippets. Using a Delphi-style protocol, multiple experienced developers iteratively ranked snippets by comprehension difficulty and resolved disagreements through structured feedback. The resulting rankings provided a practical approximation of ground-truth comprehensibility. Because expert-based evaluation is costly and does not scale, we evaluated a range of proxy measures using student studies. The results showed that proxies based on input-output reasoning about program behavior align most strongly with expert judgments. In particular, the time required to answer input-output questions was the most reliable proxy.

Based on these findings, we adopt input-output comprehension tasks and use response time as the primary measure of comprehensibility in this study. This measure captures the effort required to understand the behavior of the code while remaining scalable and suitable for controlled experiments.

\subsection{Methodology}

\subsubsection{Study Overview}

The goal of this study is to estimate the effect of code verifiability on code comprehensibility under a controlled experimental setting. We conduct a controlled experiment in which participants perform input-output comprehension tasks on semantically equivalent versions of Java methods that differ only in their verifiability with respect to the Checker Framework. Verifiability is operationalized relative to a fixed verification ecosystem, and is manipulated through semantics-preserving code transformations that induce or eliminate false positive warnings produced by selected checkers.

The study is designed as a multi-item within-subjects experiment, where each participant completes multiple tasks across different code snippets, and each code snippet appears in two versions across participants, with only one version shown to each participant. The analysis is conducted using mixed-effects models to account for variability across participants and snippets. We adopt a between-subjects design to avoid learning and carryover effects, as exposure to one version of a snippet may influence how participants interpret the other version.

\subsubsection{Independent Variable and Treatment Definition}

The primary independent variable is \textit{verifiability condition}, defined relative to a fixed verification setting based on the Checker Framework.

This variable is manipulated as a binary factor with two levels:
\begin{itemize}
    \item \textbf{High-verifiability condition}: the code snippet produces no false positive warnings under the selected checker.
    \item \textbf{Low-verifiability condition}: a semantically equivalent version of the same snippet produces one or more false positive warnings under the same checker.
\end{itemize}

The treatment is therefore not an abstract notion of verifiability, but a concrete manipulation of code that changes whether the checker can successfully verify certain safety properties. This manipulation is implemented through controlled code transformations that preserve program semantics while affecting the ability of the checker to reason about the code. 

We define semantics preservation in terms of input--output behavior: for any valid input, the transformed program produces the same observable output (including return values and externally visible side effects) as the original program. Our transformations are not required to preserve lower-level properties such as memory layout, evaluation order, or intermediate states, as long as the externally observable behavior remains unchanged.

\subsubsection{Dependent Variable}

The dependent variable is \textit{code comprehensibility}, operationalized using a validated proxy measure.

The primary outcome is the \textit{time required to answer an input-output comprehension question} for a given code snippet. This measure captures the effort required for a participant to understand the behavior of the code. Time is recorded from the moment the participant begins the task until an answer is submitted.

Correctness is recorded as a secondary measure to ensure data quality.

\subsubsection{Treatment Realization: Code Transformations}

The manipulation of verifiability is achieved through semantics-preserving code transformations that induce false positive warnings.

To illustrate the type of transformation we consider, we present a simplified example based on the index checker from the Checker Framework (Figure~\ref{fig:verifiability-example}). 

\begin{figure}[t]
\centering
\begin{minipage}{0.48\textwidth}
\textbf{Low-verifiability version}
\begin{lstlisting}[language=Java]
static boolean isSafe(int i, int[] a) {
    return i >= 1 && i < a.length;
}

static void f(int[] a, int i) {
    if (isSafe(i, a)) {
        int x = a[i]; // Warning
    }
}
\end{lstlisting}
\end{minipage}
\hfill
\begin{minipage}{0.48\textwidth}
\textbf{High-verifiability version}
\begin{lstlisting}[language=Java]
static void f(int[] a, int i) {
    if (i >= 1 && i < a.length) {
        int x = a[i]; // No warning
    }
}
\end{lstlisting}
\end{minipage}
\caption{Example of a semantics-preserving transformation that affects verifiability, based on the Index Checker of the \emph{Checker Framework}. The low-verifiability version produces false positive warnings (\lstinline{[array.access.unsafe.low]} and \lstinline{[array.access.unsafe.high]}) because the bounds check is encapsulated in a helper method, preventing the checker from refining constraints across method boundaries. As a result, it cannot establish that the array access is safe and conservatively issues warnings to preserve soundness. In contrast, the high-verifiability version performs the bounds check directly in the conditional, enabling the checker to refine the value of \lstinline{i} within the guarded block and verify that the access is safe, resulting in no warnings.}
\label{fig:verifiability-example}
\end{figure}

\subsection{Snippet Selection and Construction}
\label{sec:snippet-selection-construction}

We construct a dataset of code snippets through a multi-stage pipeline that (1) selects representative open-source projects, (2) filters candidate methods, (3) extracts minimized executable programs, and (4) constructs pairs of semantically equivalent variants that differ in verifiability.

\paragraph{Project Selection}
We begin by identifying candidate open-source repositories hosted on GitHub using repository mining tools such as SEART and CIRT. Our goal is to select projects representative of real-world, professional software development rather than hobby or toy projects, to improve the external validity of our results.

We apply the following inclusion criteria:
\begin{itemize}
    \item \textbf{Language:} The project must be written in Java. We choose Java because it is widely used in industry and is the most common language in prior code comprehensibility studies.
    \item \textbf{Maturity and popularity:} The repository must have at least 1{,}000 commits, 1{,}000 pull requests, and 1{,}000 stars. We use these thresholds as a proxy for identifying mature, high-quality, and well-maintained projects.
    \item \textbf{Recent activity:} The repository must contain at least one commit within three weeks prior to the data collection date, ensuring that it reflects current development practices.
\end{itemize}

\paragraph{Method Filtering}
From the selected repositories, we extract candidate methods that can serve as the core of our snippets. In our study, a \emph{snippet} is defined as a minimal executable program whose central logic is a target method.

We select methods satisfying the following criteria:
\begin{enumerate}
    \item All types used in the method signature and body must come exclusively from the Java Development Kit (JDK).
    \item The method must not contain annotations in its signature.
    \item The method must be declared as \texttt{static}.
    \item The method must include a JavaDoc comment.
    \item The method must have at least one parameter and a non-\texttt{void} return value.
    \item The method must contain between 20 and 30 non-blank, non-comment lines of code.
\end{enumerate}

Criteria (1)--(3) ensure that snippets are self-contained and do not require project-specific knowledge. Criterion (5) is required by our experimental design, which uses input--output comprehension tasks. Criterion (6) ensures that tasks remain within a reasonable time budget while still capturing meaningful behavior. Applying these filters yields a set of $N$ candidate methods.

\paragraph{Program Slicing}
For each candidate method, we construct a minimized executable program using \emph{Specimen}, a program slicing tool. The resulting program contains only the code necessary to preserve the behavior relevant to the target method.

The key requirement is that, for a given verification tool, the analysis result (i.e., the presence or absence of warnings) on the sliced program matches the result obtained when analyzing the full project.

\paragraph{Verifiability-Controlled Construction}
We construct snippet pairs using a two-phase procedure.

\subparagraph{Phase 1: Eliminating False Positives}
We run a selected set of Checker Framework checkers on all minimized programs. For each method that produces a warning, we manually inspect whether the warning corresponds to a real defect or a false positive.

For methods with false positive warnings, we apply small transformations intended to help the checker establish a proof and eliminate the warning. These transformations are designed to preserve program behavior (defined formally below). If a transformation successfully removes the warning without changing behavior, we retain both versions:
\begin{itemize}
    \item the original version (warning-producing): \emph{low-verifiability}
    \item the transformed version (warning-free): \emph{high-verifiability}
\end{itemize}

Phase~1 may favor methods that naturally produce false positives. We address this by including Phase~2, which systematically injects false positives into otherwise warning-free methods, improving coverage and diversity.

\subparagraph{Phase 2: Injecting False Positives}
To mitigate selection bias, we also construct candidates in the reverse direction. For each checker, we consider methods for which no warning is initially produced and attempt to introduce false positive warnings using transformations identified in Phase~1.

In this case:
\begin{itemize}
    \item the original version (no warning): \emph{high-verifiability}
    \item the transformed version (warning-producing): \emph{low-verifiability}
\end{itemize}

Transformations are applied only at suitable program locations. For example, nullness-related warnings can be introduced at dereference sites, while index-related warnings require array accesses. This constraint ensures that injected warnings are realistic and avoids introducing artificial patterns.

These transformations are not treated as independent experimental factors. Instead, they serve as different realizations of the treatment condition. They are included in the analysis to account for heterogeneity in how the treatment is implemented.

\paragraph{Semantic Preservation}
All transformations are required to preserve the functional behavior of the method. We define this using input--output equivalence.

Let $m$ and $m'$ be two versions of a method with the same signature. We say that $m'$ is \emph{semantics-preserving} with respect to $m$ if:
\[
\forall x \in \mathcal{I}, \quad m(x) = m'(x),
\]
where $\mathcal{I}$ is the set of all valid inputs to the method, and $m(x)$ denotes the output produced by executing $m$ on input $x$.

We manually validate that each transformation satisfies this property.

\paragraph{Final Dataset Construction}
The union of candidates produced in Phases~1 and~2 forms a pool of snippet pairs. Each pair consists of two semantically equivalent versions that differ only in whether the checker can verify the code.

From this pool, we construct the final dataset. Based on prior studies and estimated task duration, we target at most 20 snippets. If more than 20 candidates are available, we sample a balanced subset across phases to reduce potential confounding effects. The final selection is further validated through pilot studies to ensure feasibility within the time budget.

\subsubsection{Participants}

Participants are undergraduate computer science students with prior experience in Java programming. Eligibility criteria include completion of relevant programming courses and familiarity with basic control flow and data structures.

Participants are recruited through standard academic channels such as flyers, course announcements, and email invitations. All participants are at least 18 years old and provide informed consent prior to participation.

\subsubsection{Experimental Procedure}

The study is conducted in a single session using a web-based interface.

Each participant completes tasks on a sequence of code snippets presented one at a time. For each snippet, the participant is shown the code and asked to determine the output for a given input. The code remains visible during the task to avoid confounding memory effects.

Each participant sees only one version of each underlying snippet, either the high-verifiability or the low-verifiability version. Assignment of versions is randomized across participants so that, for each snippet, both versions are observed across the participant pool.

The order of snippets is randomized for each participant to mitigate order effects. Participants are not allowed to return to previous tasks once they proceed.

The interface records:
\begin{itemize}
    \item time spent reading the code,
    \item time spent answering the question,
    \item the final answer provided by the participant.
\end{itemize}

Participants may pause the timer when not actively working. Short breaks may be incorporated to reduce fatigue.


\subsubsection{Statistical Analysis}

We use linear mixed-effects models, which are standard for analyzing repeated-measures data with both participant- and item-level variability. This approach allows us to account for differences across participants and snippets while estimating the fixed effect of verifiability.

The primary model includes:
\begin{itemize}
    \item a fixed effect for verifiability condition,
    \item random intercepts for participants and snippets,
    \item a covariate for snippet order.
\end{itemize}

Additional fixed effects include checker type, transformation type, and measured code characteristics (e.g., code length), to account for variation across these dimensions.

The primary quantity of interest is the effect of verifiability condition on comprehension time. Statistical significance is assessed using standard inferential procedures on the model coefficients.

\subsection{Implications}
This study provides evidence about the relationship between code verifiability and code comprehensibility under controlled conditions. 

If a strong relationship is observed, it would support the use of verifiability as a practical automatic proxy for comprehensibility, enabling scalable analysis and tool support without requiring costly human studies. 

If no substantial relationship is found, the results would suggest that prior observed associations may have been influenced by confounding factors, highlighting the need for more precise definitions and measurements of both constructs.

In either case, the study contributes to a clearer understanding of the relationship between verification and human-centered software properties, and informs future research on automated proxies for code quality.

\subsection{Control of Confounding Factors}
\label{sec:confounder-control}

The study design systematically addresses the confounding factors identified in Section~\ref{sec:confounders}. We incorporate controls at the participant, snippet, and experimental levels to isolate the effect of verifiability.

\textbf{Participant-level confounders.}
Differences in programming experience, domain knowledge, and general ability are partially controlled through eligibility criteria that ensure a relatively homogeneous population of undergraduate students with prior Java experience. Remaining variability is accounted for statistically by including participant as a random effect, and by collecting background measures (e.g., self-reported experience) that are included as covariates.

\textbf{Snippet-level confounders.}
Variation in inherent difficulty, code structure, and familiarity is controlled through a paired design in which each underlying method appears in both high- and low-verifiability conditions across participants. This ensures that comparisons are made between semantically equivalent versions of the same code. Transformations are selected to preserve semantics while minimizing changes to code length, control flow, and layout. Residual structural variation (e.g., code length) is measured and included as a covariate.

\textbf{Experimental confounders.}
Learning effects, fatigue, and ordering effects are mitigated by randomizing the order of snippet presentation for each participant and limiting session length. Snippet order is explicitly modeled as a covariate. Instrumentation-related confounders are reduced by using a standardized web-based interface with consistent task presentation, timing mechanisms, and instructions.

\textbf{Verification-context confounders.}
To ensure a well-defined notion of verifiability, the verification context is fixed by selecting a specific set of Checker Framework checkers and configurations. This prevents variation in tool behavior from influencing the results.

Together, these design choices ensure that the study directly addresses the confounders identified in Section~\ref{sec:confounders}, enabling a controlled evaluation of the effect of verifiability on comprehensibility.

\subsection{Threats to Validity}
\label{sec:threats}

\textbf{Internal validity.}
The study aims to isolate the effect of verifiability by using semantically equivalent code snippets that differ only in whether the checker produces false positive warnings. However, semantics-preserving transformations may introduce unintended structural differences (e.g., minor changes in control flow or code layout) that could influence comprehension. We mitigate this risk by restricting transformations to simple, well-understood patterns and by manually validating input--output equivalence. Residual structural differences (e.g., code length) are measured and included as covariates in the analysis. 

Another potential threat arises from how false positive warnings are identified. Misclassification of warnings (i.e., treating a true defect as a false positive or vice versa) could affect the validity of the manipulation. To mitigate this risk, warnings are inspected carefully, and decisions are informed by expertise with the Checker Framework.

\textbf{Construct validity.}
Verifiability is operationalized as the presence or absence of false positive warnings produced by a fixed set of Checker Framework checkers. While this captures an important aspect of verification effort, it does not fully represent all dimensions of verifiability, such as the effort required to write specifications or interact with verification tools.

Code comprehensibility is measured using response time on input--output comprehension tasks. Although prior work has shown that this measure aligns with expert judgments, it remains an indirect proxy and may not capture all aspects of human understanding.

\textbf{External validity.}
The study focuses on Java methods drawn from mature open-source projects and evaluates verifiability using a specific set of Checker Framework checkers. As a result, the findings may not generalize to other programming languages, domains, or verification tools. In addition, participants are undergraduate students, which may limit generalizability to professional developers.

\textbf{Statistical conclusion validity.}
The study relies on linear mixed-effects models to estimate the effect of verifiability on comprehension time. Threats include insufficient statistical power, noise in response time measurements, and potential model misspecification. We mitigate these risks by using repeated measures across multiple snippets and participants, and by including random effects for participants and snippets as well as relevant covariates.


